% Part: sets-functions-relations
% Chapter: sets
% Section: important-sets

\documentclass[../../../include/open-logic-section]{subfiles}

\begin{document}

\olfileid{sfr}{set}{imp}
\olsection{Enkele belangrijke verzamelingen}

\begin{ex}
We zullen voornamelijk verzamelingen behandelen waarvan de !!{element}en
wiskundige objecten zijn. Vier zulke verzamelingen zijn belangrijk genoeg
om een eigen naam te hebben:
\begin{multline*}
    \Nat = \{0, 1, 2, 3, \ldots\} \\
    \shoveright{\text{de verzameling natuurlijke getallen}}\\
    \shoveleft{\Int = \{\ldots, -2, -1, 0, 1, 2, \ldots\}} \\
    \shoveright{\text{de verzameling gehele getallen}}\\
    \shoveleft{\Rat = \Setabs{\nicefrac{m}{n}}{m, n \in \Int\text{ en }n \neq 0}}\\
    \shoveright{\text{de verzameling rationale getallen}}\\
    \shoveleft{\Real = (-\infty, \infty)}\\
    \text{de verzameling reële getallen (het continuüm)}
\end{multline*}
Dit zijn allemaal \emph{oneindige} verzamelingen: elk van deze verzamelingen
heeft oneindig veel !!{element}en.

Wanneer we deze verzamelingen in deze volgorde doorlopen, voegen we telkens
\emph{meer} getallen toe. De inclusies $\Nat \subseteq \Int
\subseteq \Rat \subseteq \Real$ liggen voor de hand: ieder natuurlijk
getal is immers een geheel getal, ieder geheel getal is rationaal en ieder
rationaal getal is reëel. Ook $\Nat \subsetneq \Int \subsetneq
\Rat$ is duidelijk, want $-1$ is wel een geheel maar geen natuurlijk getal,
en $\nicefrac{1}{2}$ is wel rationaal maar niet geheel. Minder duidelijk is
dat $\Rat \subsetneq \Real$, dus dat er reële getallen bestaan die niet
rationaal zijn\oliflabeldef{sfr:arith:real:realline}{; hierop komen we terug
in \olref[arith][real]{realline}}{}.

Soms gebruiken we ook de verzameling positieve gehele getallen $\PosInt = \{1,
2, 3, \dots\}$ en de verzameling die alleen de eerste twee natuurlijke
getallen bevat, $\Bin = \{0, 1\}$.
\end{ex}

\begin{tagblock}{compsci}
\begin{ex}[Tekenreeksen]
Een ander interessant voorbeeld is de verzameling $A^{*}$ van \emph{eindige
tekenreeksen} over een alfabet $A$: iedere eindige rij elementen van~$A$
is een tekenreeks over $A$. De \emph{lege tekenreeks $\Lambda$} rekenen we
tot de tekenreeksen over~$A$, voor elk alfabet~$A$. Bijvoorbeeld:
\begin{multline*}
\Bin^*
=\{\Lambda,0,1,00,01,10,11,\\
000,001,010,011,100,101,110,111,0000,\ldots\}.
\end{multline*}
Als $x=x_{1}\ldots x_{n}\in A^{*}$ een tekenreeks is die uit $n$
``letters'' van $A$ bestaat, dan heet~$n$ de \emph{lengte} van de tekenreeks
en schrijven we $\len{x}=n$.
\end{ex}
\end{tagblock}

\begin{ex}[Oneindige rijen]
Voor elke verzameling $A$ kunnen we ook de verzameling~$A^\omega$ van
oneindige rijen !!{element}en van~$A$ beschouwen. Een oneindige rij
$a_1a_2a_3a_4\dots$ bestaat uit een eenzijdig oneindige lijst van objecten,
die elk !!a{element} van~$A$ zijn.
\end{ex}
\end{document}
